\section{The problem, the results, and the mechanisms}\label{sec:question}
Let $X=(X_1,\ldots,X_n)$ be uniform on $\{-1,1\}^n$, and obtain $Y$ by independently reversing each sign with probability $p\in[0,1/2]$; put $\rho=1-2p$. A deterministic Boolean function $f:\{-1,1\}^n\to\{0,1\}$ retains one bit of $X$. The question is whether combining coordinates can retain more information about $Y$ than simply keeping one coordinate.

All logarithms are natural. Write
\[
 h(s)=-s\ln s-(1-s)\ln(1-s),\qquad 0\ln0:=0,
\]
and let $\mu=\E f(X)$. Then
\begin{equation}\label{eq:information}
\II(f(X);Y)=h(\mu)-\HH(f(X)\mid Y).
\end{equation}
For $f(x)=\one\{x_i=1\}$, the right side is $\ln2-h(p)$. Courtade and Kumar conjectured that this is the optimum over all Boolean functions~\cite{CK}.

\begin{theorem}[Courtade--Kumar inequality]\label{thm:main}
For every finite $n$, every Boolean $f$, and every $p\in[0,1/2]$,
\begin{equation}\label{eq:target}
\II(f(X);Y)\le\ln2-h(p).
\end{equation}
\end{theorem}
The statement uses uniform independent input bits and independent binary symmetric noise. It does not concern a general input distribution or a general channel. Expressing every entropy in bits divides both sides by $\ln2$.

\subsection{Quantitative content}
The proof gives more than an extremal value. At low noise it establishes the mean-dependent profile
\begin{equation}\label{eq:profile}
\HH(f(X)\mid Y)\ge V(\mu)h(p),\qquad V(\mu)=4\mu(1-\mu).
\end{equation}
This holds for $0\le p\le1/20$. The factor $V(\mu)$ is four times the output variance. It is the quantity that makes restriction into smaller Boolean functions compatible with induction.

For $n\ge1$, define the distance to a coordinate, allowing complementation, by
\begin{equation}\label{eq:delta-intro}
 \delta(f)=\min_{i,\,\sigma\in\{-1,1\}}\Pp\{2f(X)-1\ne\sigma X_i\}.
\end{equation}
The quantitative conclusion is
\begin{equation}\label{eq:stability-overview}
\ln2-h(p)-\II(f(X);Y)
\ge \frac{p(1-2p)^2}{2000}\,\delta(f),
\qquad 0\le p\le\frac12.
\end{equation}
Thus, at a fixed $0<p<1/2$, near-optimal information forces closeness to a coordinate, independently of $n$. The explicit constants are sufficient, not sharp. Section~\ref{sec:use} gives a stronger piecewise coefficient, a matching qualitative converse, and worked examples. The stability results are consequences of the argument, not premises of Theorem~\ref{thm:main}.

\subsection{Three ideas that can be used separately}
\paragraph{Restriction needs a compensating quantity.}
If the two coordinate sections have means $\mu_0,\mu_1$ and $d=|\mu_1-\mu_0|$, then
\[
 \frac{V(\mu_0)+V(\mu_1)}2=V(\mu)-d^2.
\]
A local entropy inequality must pay for this precise loss. Proposition~\ref{prop:induction} formulates the resulting scalar-to-cube implication with a free profile coefficient $c$; it is not tied to the numerical cutoff used in the application.

\paragraph{The useful symmetric quantity is a product.}
For two posterior probabilities, let $J$ be the average entropy before mixing, $M$ the average entropy after mixing, and $K=M-J$ the gain. The scalar problem has the form
\[
 K+Md^2-2\ell d
 =M\left(d-\frac{\ell}{M}\right)^2+\frac{KM-\ell^2}{M}.
\]
This is why the product $KM$ is compared with its symmetric value. Lemma~\ref{lem:product} proves that both $K$ and $KM$ are smallest at symmetry for a fixed separation. Its noise range is the entire open interval $0<p<1/2$, although the stronger two-point estimate needed for CK is used only at low noise.

\paragraph{A scalar envelope turns spectral information into entropy.}
Away from low noise, a bound on constant and linear Fourier mass controls a quadratic statistic of the posterior. The quadratic is dictated by a spectral cancellation, not guessed. Proposition~\ref{prop:certificate} states the resulting information certificate with a free Fourier cap and a free entropy envelope. It separates what is special to the application from what can be reused.

\subsection{Proof organization}
Sections~\ref{sec:induction}--\ref{sec:product} establish the low-noise branch. Sections~\ref{sec:threeclasses}--\ref{sec:energy} establish the complementary branch, and Section~\ref{sec:completion} joins them at $p=1/20$. Section~\ref{app:linear} quantifies the slack. Section~\ref{sec:use} explains how to apply the conclusions and what their constants do not say. Appendices~\ref{app:centered}--\ref{app:constants} contain every deferred scalar and finite-constant argument.

For the main theorem, the task is therefore finite and explicit: prove one all-posterior mixing statement at low noise, and one Fourier-to-entropy estimate for the complementary class. No finite truth-table experiment is used to infer a statement in arbitrary dimension.

\subsection{Relation to earlier work}
Samorodnitsky proved the Courtade--Kumar inequality in a dimension-independent high-noise interval~\cite{Sam}. Yu's extended paper proves local optimality among balanced functions and a computer-assisted balanced result for $0\le\rho\le0.914$~\cite{YuLocal}. Javanmard and Woodruff prove an arbitrary-bias coordinate-wise information inequality and sharpen high-noise entropy estimates~\cite{JW}. Their coordinate-wise sum is a different quantity from information about the full noisy vector.

The present proposed argument addresses all output means, all finite dimensions, and all noise levels, with the explicit stability estimate above. The ordinary variance identity, antipodal projection, and uniform-reference contraction are not claimed as new. First-level stability is classical~\cite{FKN}. The mathematical content to examine here is the compatibility of the local mixing inequality with restriction, the entropy-product comparison, and their combination with the Fourier certificate. None of the cited specialized results is a premise of the proof.

\subsection{Notation and normalization}
\begin{equation}\label{eq:eta}
\rho=1-2p,\qquad
\eta(u)=h\!\left(\frac{1+u}{2}\right),\qquad
\kappa(u)=\ln2-\eta(u),\quad -1\le u\le1.
\end{equation}
Thus $\eta(\rho)=h(p)$ and $\kappa(\rho)=\ln2-h(p)$. The indicator $f$ and the sign function $F=2f-1$ encode the same output bit; write $m=\E F=2\mu-1$. A signed coordinate means $F(x)=\sigma x_i$, or equivalently $f(x)=\one\{\sigma x_i=1\}$. Fourier weights are introduced in Section~\ref{sec:geometric}.

On $(-1,1)$,
\[
 \eta'(u)=-\atanh u,\qquad \eta''(u)=-\frac1{1-u^2}.
\]
Integrating the power series for $\atanh$ gives
\begin{equation}\label{eq:series}
\kappa(u)=\sum_{j\ge1}\frac{u^{2j}}{(2j)(2j-1)}.
\end{equation}
The coefficients are positive, and hence
\begin{equation}\label{eq:entropyfacts}
\kappa(u)\ge\frac{u^2}{2},\qquad
\kappa(\theta u)\le\theta^2\kappa(u)\quad(0\le\theta\le1).
\end{equation}
Endpoint values follow by continuity. The numerical comparisons used below are proved by rational enclosures in Appendix~\ref{app:constants}.

Two parameters must remain distinct: $d$ is the fixed difference of section means; $t=|a-b|$ is the separation of a posterior pair. In an application $|A-B|$ varies with the noisy observation while $d$ remains fixed. No independence between $A$ and $B$ is required.
